\chapter{Classical Mechanics}
\label{chap:classical-mechanics}

% Closed question: What selects one admissible continuation as the history
% actually carried forward?

\begin{chapteressay}

A ball rolls off a table.

At one moment it is at the table edge. A little later it is lower and farther
out. A camera can mark those two positions. A stopwatch can mark the elapsed
time. Between the two marks, no one sees every infinitesimal point of the
motion. Still, the path is not arbitrary. The ball does not zigzag, pause,
circle, and resume. It follows the familiar parabola.

That parabola is the chapter's first example of an admissible continuation. The
measured endpoints do not contain the whole motion as a list. They contain
constraints. Classical mechanics supplies a selection rule: among the possible
histories compatible with the endpoints, the actual history is the one whose
action is stationary. The law is often presented as a postulate. In this book
it is read as a measurement consequence. A record with endpoints needs a way to
choose a continuation that does not add structure the measurement never earned.

The speedometer has been doing a low-tech version of this all along. Between
two pulses, it assumes the least fussy continuation compatible with the last
known state. It does not imagine that the wheel sped up, slowed down, reversed,
and recovered unless the later record forces that complication. The odometer
accumulates distance by trusting that local continuation. If the next pulse
disagrees, the interpolant is corrected. If the pulses remain stable, the
minimal continuation survives.

The same structure appears in astronomy. Mars is observed at one opposition and
then again months later. The observations are finite. The orbit is not drawn
on the sky as a luminous thread. Yet the inverse-square law, the Sun's position,
and the elapsed time select an ellipse. A wild path between the two observed
locations may be imaginable, but it would require forces not present in the
record. It would be an added story.

The code names this danger with comic severity. \lean{Spline} has observations,
knots, and interpolants. \lean{GalerkinProcess} supplies a polynomial carrier
for turning finite evidence into admissible approximation. \lean{ArmWaveProcess}
can reticulate a guess into a more elaborate spline. \lean{BULLSHIT} then asks
whether the interpolation really respects the record. The joke is doing
technical work: overhead is conserved. A story that adds unmeasured structure
must carry the burden of that addition.

The closed question is:

\begin{center}
\emph{What selects one admissible continuation as the history actually carried
forward?}
\end{center}

The answer is not that nature loves parabolas, ellipses, or smooth curves as
decorations. The answer is that a physical history is the least added structure
compatible with the measurements and the governing constraint. Least action is
the classical expression of that rule. It turns endpoint records into a
selected path.

Chapter 6 therefore gives distance a history. Chapter 5 counted refinements
between marks. Now those marks are joined. The joining is not free invention.
It must be accountable to the record.

\end{chapteressay}

% Phenomenon arcs use: 1/3 setup -> phenomenon box -> 2/3 structural interpretation.

\section{Least Action}
\label{se:least-action}

The principle of least action can sound mystical if it is heard as a claim that
the particle surveys every possible future and chooses the best one. It should
be heard instead as a selection rule for histories.

A classical system has a Lagrangian $L$, usually kinetic energy minus potential
energy. A proposed path $q(t)$ receives an action

\[
S[q] = \int_{t_0}^{t_1} L(q,\dot q,t)\,dt.
\label{eq:chapter6-action}
\]

The physical path is stationary under small variations that leave the endpoints
fixed. For a ball in a uniform gravitational field, that condition gives a
parabola. For a free particle, it gives a straight line at constant velocity.
For a pendulum at small angle, it gives harmonic motion.

In measurement language, the endpoints are the record, and the action is the
burden paid by a proposed continuation. A path that introduces needless bends,
accelerations, or reversals pays extra burden without evidence. Least action
does not say that nothing happens between marks. It says that the history
between marks must be justified by the constraint actually present.

This is where the speedometer analogy becomes precise. A constant wheel speed
between two pulses is not a metaphysical truth. It is the minimal admissible
interpolant until the next pulse forces a revision. If the pulses begin to
arrive closer together, acceleration enters the record. If the pulses spread
out, deceleration enters. Until then, the interpolant should not add drama.

The Lean code's \lean{Spline} type is a useful grammar for this. An
\lean{observation} starts the record. A \lean{knot} binds a proposition to a
polynomial carrier. An \lean{interpolant} combines facts and polynomials across
prior spline structure. Classical mechanics is not being implemented there, but
the grammar is exactly the one this chapter needs: finite marks, admissible
carriers, and composed continuations.

\section{Euler-Lagrange from Data}
\label{se:euler-lagrange-from-data}

The Euler-Lagrange equations are what the principle of least action looks like
when the stationary condition is written locally. For a Lagrangian $L(q,
\dot{q}, t)$ describing a system with generalized coordinate $q(t)$, the
equation is
\[
\frac{d}{dt}\frac{\partial L}{\partial \dot q}
  - \frac{\partial L}{\partial q} = 0.
\label{eq:chapter6-euler-lagrange}
\]
This is a second-order ordinary differential equation for $q(t)$. Given
appropriate boundary conditions --- typically the position at two endpoint
times --- the equation has a unique solution (under mild regularity
assumptions on $L$). The equation is not a replacement for data. It is the
rule that tells a data record what continuations remain admissible.

The derivation is brief but worth retracing because the structure carries
the chapter's argument. The action functional is
\[
S[q] = \int_{t_1}^{t_2} L(q(t), \dot{q}(t), t) \, dt,
\]
where $q(t_1)$ and $q(t_2)$ are fixed boundary data. Consider a small
variation $q \to q + \varepsilon \eta$ with $\eta(t_1) = \eta(t_2) = 0$.
The first-order change in $S$ is
\[
\delta S = \int_{t_1}^{t_2}\left(\frac{\partial L}{\partial q} \eta + \frac{\partial L}{\partial \dot{q}} \dot{\eta}\right) dt.
\]
Integrating the second term by parts (using $\eta(t_1) = \eta(t_2) = 0$ to
kill the boundary):
\[
\delta S = \int_{t_1}^{t_2}\left(\frac{\partial L}{\partial q} - \frac{d}{dt}\frac{\partial L}{\partial \dot{q}}\right) \eta \, dt.
\]
For $S$ to be stationary --- $\delta S = 0$ for all admissible $\eta$ --- the
integrand's coefficient of $\eta$ must vanish, which gives the
Euler-Lagrange equation. The derivation is local: at every interior point
between $t_1$ and $t_2$, the equation must hold.

Now apply this to a planet. Suppose Mars is observed at two specific
positions $\mathbf{r}_1$ and $\mathbf{r}_2$ separated by six months of
proper time. The Lagrangian for a planet of mass $m$ in the Sun's
gravitational field (with the Sun fixed at the origin) is
\[
L = \frac{1}{2} m |\dot{\mathbf{r}}|^2 + \frac{GMm}{|\mathbf{r}|},
\]
where $G$ is Newton's constant and $M$ is the Sun's mass. The Euler-Lagrange
equation for this Lagrangian is
\[
m\ddot{\mathbf{r}} = -\frac{GMm}{|\mathbf{r}|^3} \mathbf{r}.
\]
This is Newton's equation of motion for an inverse-square gravitational
force directed toward the Sun. Given the boundary data $\mathbf{r}(t_1) =
\mathbf{r}_1$ and $\mathbf{r}(t_2) = \mathbf{r}_2$, the equation has a
unique solution (modulo the orbital plane orientation, which is determined
by the angular momentum at any one point).

The solution is a conic section: ellipse, parabola, or hyperbola
depending on the energy. For a bound planet like Mars, the orbit is an
ellipse with the Sun at one focus. The specific ellipse is fixed by the
boundary data: given two positions separated by a known time interval,
the ellipse's semi-major axis, eccentricity, orientation, and orbital
phase are all determined.

The point is that the orbit is not chosen; it is forced. Among all
admissible continuations of the data (all curves passing through
$\mathbf{r}_1$ at $t_1$ and $\mathbf{r}_2$ at $t_2$), the unique one
that satisfies Newton's law (equivalently, the Euler-Lagrange equation
of the gravitational Lagrangian) is the specific Keplerian ellipse.
Any other continuation would require unrecorded forces between the two
boundary times. The Euler-Lagrange equation rejects continuations that
would need such forces.

This is the chapter's central structural claim. Classical mechanics is
sometimes taught as if the equations of motion float above measurement
and are later applied. The framework reverses this. A record of
measurements asks for a continuation between two recorded events. The
Euler-Lagrange equation is the local condition that closes the question
by rejecting any continuation that would require force not present in the
recorded sources.

Without observation data --- without specifying boundary conditions ---
the Euler-Lagrange equation has infinitely many solutions parameterized
by the choice of initial conditions. The equation alone does not produce
a unique trajectory; it characterizes a family. The record's role is to
select one member of the family by specifying enough boundary data.

The \lean{GalerkinProcess} name from the Lean code is apt here. In
numerical analysis, a Galerkin method chooses an approximate solution
from a finite-dimensional space of basis functions and demands that the
residual vanish against an admissible test space. The procedure converts
a continuous problem (the Euler-Lagrange equation) into a finite linear
system, solvable on a computer. The continuous solution emerges as the
limit of finite Galerkin approximations as the basis is refined. The
chapter's claim is that this is what measurement is doing: a finite
record of boundary data plus a Lagrangian produces a finite-dimensional
admissible continuation; finer records produce finer continuations; the
limit is the continuous trajectory.

The phrase ``from data'' matters. The Euler-Lagrange equation is not a
metaphysical truth that exists independently of measurement. It is a
specific local condition that closes the gap between recorded events.
Without records, the equation is vacuous; with records, the equation
selects a unique continuation among admissible alternatives. The
discipline of measurement is to provide the records; the discipline of
mechanics is to apply the equation.

This also explains why a bad model can fit two endpoints but fail as
mechanics. A polynomial curve can be drawn through any two points. A
physical trajectory must carry the right burden: the inverse-square
force, the specific Lagrangian, the conservation laws that follow from
the action's symmetries. If the proposed curve needs a hidden impulse
in the middle, the record must show that impulse. Otherwise the curve
is just overhead --- unrecorded structure imposed on the data without
empirical warrant.

The chapter's framework formalizes this in the typeclass cascade. A
\lean{GalerkinProcess} accepts boundary data and a basis; it produces the
finite-dimensional approximation; the \lean{FINITE\_ELEPHANT} certificate
confirms that further refinement would produce no admissible distinction
at the stated resolution; the resulting trajectory is the unique
continuation forced by the data, the law, and the resolution. None of
these elements is optional. Removing any of them allows continuations
the chapter considers inadmissible.

\section{Repeatability of Invisible Motion}
\label{se:repeatability-of-invisible-motion}

Galileo's inclined plane gives the central phenomenon. A ball rolls down a
shallow ramp. The motion is too fast to measure comfortably if the ball falls
freely, so the ramp slows the descent while preserving a simple acceleration
law. Different observers can time the ball, mark positions, and recover the
same structure.

\begin{phenom}{Repeatability of Invisible Motion}
\label{ph:repeatability-of-invisible-motion}

\PhStatement A motion that is not continuously observed can still be physically
repeatable when independent finite records select the same admissible
continuation.

\PhOrigin Galileo's inclined-plane experiments made acceleration measurable by
slowing falling motion into a repeatable laboratory process.

\PhObservation Different timing devices and repeated trials recover the same
quadratic distance-time relation within experimental error, even though the
intermediate continuum of positions is not directly listed in the record.

\PhConstraint The experiment may infer invisible intermediate motion only when
the same continuation is forced by repeatable endpoint and timing records.

\PhConsequence Classical history is not guessed between observations. It is
selected by the continuation that repeatable finite measurements keep forcing.

\end{phenom}

The phenomenon closes a common loophole. If motion between observations were
mere storytelling, repeatability would not be enough to stabilize it. But
Galileo's inclined plane makes the invisible interval accountable. The ball's
distance grows quadratically with time. Change the clock, repeat the trial,
mark a different set of intermediate positions, and the same continuation
returns.

This is exactly why \lean{BULLSHIT} is a useful, if impolite, name for the
class stack. An interpolation that cannot be made repeatable is not merely
uncertain; it has imported material that the measurement cannot carry. The
burden has to be paid somewhere. Classical mechanics pays it with forces,
constraints, and action.

\section{The Free Parameter}
\label{se:the-free-parameter}

Two endpoint positions do not always determine a unique history. A thrown ball
can pass through the same two points with different launch speeds if the time
of flight is not fixed. Add the elapsed time, and the family narrows. Add a
launch angle or initial velocity, and it closes further.

The free parameter is not a failure of measurement. It is the honest shape of
the record. The record says exactly how much it knows.

In spline language, an interpolant often has degrees of freedom left after it
passes through the knots. Smoothness may remove some. Boundary derivatives may
remove more. A natural cubic spline, for instance, uses endpoint conditions to
settle curvature that the interior knots alone do not determine. The remaining
choice is not mystical; it is the place where another measurement can enter.

Classical mechanics treats the initial velocity in the same way. Position at
one time is not a whole state for a second-order law. A second-order law needs
position and velocity, or enough endpoint information to determine both. The
missing velocity is not hidden ontology. It is an unresolved parameter in the
measurement ledger.

This helps explain why the chapter does not claim that least action removes all
ambiguity by itself. Least action selects among paths with the same boundary
data. If the boundary data are incomplete, the selection remains a family. The
right response is not to invent a path. The right response is to ask which
additional measurement closes the family.

The code's class stack embodies this discipline. \lean{FINITE\_ELEPHANT} carries
the finite matter of approximation. \lean{ArmWaveProcess} can add a guess.
\lean{BULLSHIT} tests interpolation. \lean{PROPAGANDA} and \lean{ACOLYTE} show
what happens when a story becomes institutional. The mathematical joke is that
every added layer must still answer to the record.

\section{Kepler's Orbits}
\label{se:keplers-orbits}

Kepler's orbits are the chapter's mature example. Johannes Kepler, working
with Tycho Brahe's observations of Mars from 1576 to 1601, spent about
two decades analyzing the data. The observations were sparse by modern
standards: positions noted to about one arcminute precision, distributed
across multiple oppositions of Mars over twenty-five years. From this
sparse record, Kepler extracted three laws.

The first law, published in 1609, states that planetary orbits are ellipses
with the Sun at one focus. This replaced the centuries-old Ptolemaic
assumption of circular motion (modified by epicycles to fit the data). The
ellipse fit Brahe's data far more cleanly than any combination of circles.

The second law, also from 1609, states that the line connecting a planet
to the Sun sweeps out equal areas in equal times. This is the area
conservation law; it implies that planets move faster when closer to the
Sun and slower when farther.

The third law, published in 1619, states that the square of a planet's
orbital period is proportional to the cube of its orbital semi-major axis:
\[
T^2 \propto a^3.
\label{eq:chapter6-kepler-third}
\]
Specifically, $T^2 = (4\pi^2/GM) a^3$ in modern notation, where $G$ is
Newton's constant and $M$ is the Sun's mass. The proportionality constant
is the same for all planets orbiting a given star.

Kepler discovered these laws empirically. He had no theory of gravity; he
fit curves to Brahe's data and identified the patterns. The third law in
particular was the result of decades of arithmetic: Kepler tried many
relationships between period and orbit size before he found one that worked
for all six known planets.

Newton, in 1687, derived all three laws from the inverse-square law of
universal gravitation. The derivation makes the chapter's structural point
explicit. Given the inverse-square gravitational force, the resulting
two-body equation of motion is
\[
\ddot{\mathbf{r}} = -\frac{GM}{|\mathbf{r}|^3} \mathbf{r}.
\]
For a bound orbit (negative total energy), the solution is a conic section
with the central body at one focus --- specifically, an ellipse. The
trajectory satisfies all three Kepler laws automatically: the ellipse with
Sun at focus is the first law; angular momentum conservation gives the
second law (equal areas in equal times); and the period-size relation
gives the third law.

The third law's derivation is particularly clean. For a circular orbit of
radius $r$, the gravitational acceleration $GM/r^2$ must equal the
centripetal acceleration $v^2/r = (2\pi r/T)^2/r$. Setting these equal
gives
\[
\frac{GM}{r^2} = \frac{4\pi^2 r}{T^2},
\]
which rearranges to $T^2 = (4\pi^2/GM) r^3$. The generalization to
elliptical orbits replaces the radius with the semi-major axis $a$ and
gives Kepler's third law. The proportionality constant $4\pi^2/GM$
depends only on the central body's mass, so all planets orbiting the
same star satisfy the same period-size relation.

The structural significance is the chapter's claim that the orbit is
forced, not chosen. Brahe's observations of Mars provided sparse boundary
data: positions at various dates over twenty-five years. Newton's
gravitational law provided the structural constraint: the equation of
motion. Together, the data and the law force the orbit. Kepler's three
laws are not three independent assumptions; they are three observable
features of the unique admissible continuation.

\begin{phenom}{Kepler Effect}
\label{ph:kepler-effect}

\PhStatement A sparse astronomical record can select a stable orbit when the
governing constraint makes one continuation carry less unmeasured structure
than its rivals.

\PhOrigin Kepler's analysis of planetary observations replaced circular
preference with elliptical orbits, later explained by Newtonian gravitation.

\PhObservation Planetary positions, elapsed times, and the Sun's focus agree
with ellipses and with the period-size relation $T^2 \propto a^3$.

\PhConstraint The orbit may not be chosen for geometric beauty alone. It must
be the continuation that survives the observations and the force law.

\PhConsequence Classical astronomy turns finite sky records into a physical
history by selecting the least added continuation compatible with gravitation.

\end{phenom}

Kepler is important because the record is so obviously partial. No one sees
the entire orbit at once. No one measures every point. The orbit is a
continuation that becomes trustworthy because later observations keep
agreeing with it.

The third law's confirmation is particularly clean. Kepler had data for
six planets: Mercury, Venus, Earth, Mars, Jupiter, Saturn. The ratio
$T^2/a^3$ comes out to a single number across all six planets, within
the precision of Brahe's measurements. Modern measurements with
spacecraft and high-precision astronomy extend the agreement to all
known solar-system bodies (planets, moons, asteroids, Kuiper belt
objects, comets) to many decimal places. The same ratio also applies
to exoplanets around other stars (with the appropriate stellar mass
substituted): every exoplanet detected by transit or radial velocity
methods obeys Kepler's third law with high precision. The law is one
of the most-confirmed predictions in physics.

The third law is also load-bearing for the chapter's claim that the
orbit is forced rather than fit. If Kepler's third law were merely a
phenomenological fit, it would be expected to hold approximately within
the range of Brahe's data and to diverge outside that range. Instead, the
law holds with full precision across orbits ranging from Mercury's
$88$-day period to Pluto's $248$-year period, and into the entirely
different scale of moons orbiting planets (where the proportionality
constant changes to reflect the planet's mass rather than the Sun's).
The structural form is universal; the specific proportionality constant
is local; the universality of the form is what supports the chapter's
claim of a forced continuation.

This brings the chapter back to the speedometer. The odometer is a tiny
Kepler: finite marks, a law of continuation, and repeated correction. The
planet is a grand odometer: finite observations, a gravitational law, and
a history whose burden is too low to ignore. In both cases, the record is
sparse, the continuation is forced by the law, and the agreement between
prediction and subsequent observation is the chapter's calibration check.

\begin{bridge}

The chapter's question was what selects one admissible continuation as the
history actually carried forward.

The answer is least added structure under a governing constraint. The falling
ball follows the stationary-action path. The Galilean ramp makes invisible
motion repeatable. The free parameter marks exactly what the record has not
closed. Kepler's orbit shows finite observations becoming a stable history.

The next question is sharper. Classical mechanics joins measurements by a
single admissible path. Quantum mechanics asks what happens when more than one
history remains admissible until the act of witnessing changes the record.

\end{bridge}

\begin{coda}{A Figure Skater Tracing One Clean Arc}

Two chalk marks sit on the ice. The skater begins at the first and must finish
at the second. The rink is open. Nothing prevents a spiral, a stumble, a loop,
or a sequence of decorative cuts. Many paths are imaginable.

But the skater is moving with a blade on ice.

The blade has an edge. The ice has friction. The skater has momentum. A clean
arc is not merely prettier than the alternatives. It is the path that carries
the least unnecessary correction from the given entry speed to the required
exit mark. A loop would demand extra force. A stutter would lose momentum. A
sharp corner would require a cut the blade cannot hide. The good path is the
one that makes the constraints visible without adding a story.

The coach can see this at once. If the arc is clean, the body, blade, and ice
agree. If it wobbles, the wobble is a record: the skater changed edge, lost
balance, or corrected too late. The ice preserves the evidence. The trace is a
ledger.

The analogy performs the chapter's typeclasses. The chalk marks are
observations. The blade is the carrier. The rink supplies admissibility. The
skater's body carries magnitude and load. The clean arc is the spline. The
coach's judgment is the interpolation test. A flourish that was not required by
the marks must pay for itself in visible force.

This is not about beauty, though beauty appears when burden is low. The clean
arc looks inevitable because the hidden alternatives have been disqualified by
the physical situation. The skater did not choose from fantasy paths. The
skater carried a body through a constrained medium.

Classical mechanics does the same. A path is not accepted because it can be
drawn. It is accepted because the required accelerations, forces, endpoints,
and times agree. The trace on the ice is a human-scale least-action record.

The skate mark remains after the motion. It says: this is how the continuation
was carried.

\end{coda}
